\documentclass[11pt,a4paper]{article}

% ---------------------------------------------------------------
% Packages
% ---------------------------------------------------------------
\usepackage[T1]{fontenc}
\usepackage{amsmath, amsthm, mathtools}
\usepackage{libertinus}
\usepackage{libertinust1math}
\usepackage{bm}
\usepackage{graphicx}
\usepackage{xcolor}
\usepackage{tikz}
\usetikzlibrary{arrows.meta, patterns}
\usepackage[most]{tcolorbox}
\usepackage{siunitx}
\sisetup{group-separator={\,}, group-minimum-digits=4, exponent-product=\cdot}
\DeclareSIUnit\angstrom{\text{\AA}}
\DeclareSIUnit\bohr{\ensuremath{a_0}}
\usepackage{booktabs}
\usepackage{tabularx}
\usepackage[numbers,sort&compress]{natbib}
\usepackage{microtype}
\usepackage{parskip}
\usepackage[margin=2.5cm, headheight=20pt]{geometry}
\usepackage{titlesec}
\usepackage{fancyhdr}
\usepackage{lastpage}
\usepackage{hyperref}
\usepackage[capitalise,noabbrev]{cleveref}

\numberwithin{equation}{section}

% ---------------------------------------------------------------
% Corporate colours (sampled from the FAU logo) and layout
% ---------------------------------------------------------------
\definecolor{faublue}{RGB}{0,47,108}
\definecolor{fautech}{RGB}{119,159,181}

\hypersetup{colorlinks=true, linkcolor=faublue, citecolor=faublue, urlcolor=faublue,
  pdftitle={Kinetic Monte Carlo Simulation of Diffusion},
  pdfauthor={Ali Ghelichkhani}}

\titleformat{\section}{\sffamily\Large\bfseries\color{faublue}}{\thesection}{0.75em}{}
\titleformat{\subsection}{\sffamily\large\bfseries\color{faublue}}{\thesubsection}{0.7em}{}
\titleformat{\paragraph}[runin]{\sffamily\bfseries\color{faublue}}{}{0pt}{}
\titlespacing*{\paragraph}{0pt}{1.2ex plus 0.4ex}{0.6em}

\pagestyle{fancy}
\fancyhf{}
\fancyhead[L]{\sffamily\small\color{fautech}Materials Simulation Practical~6 \textbullet{} Part~I}
\fancyhead[R]{\raisebox{-0.2\height}{\includegraphics[height=0.42cm]{figures/ww8.png}}}
\fancyfoot[L]{\sffamily\small\color{fautech}A.~Ghelichkhani}
\fancyfoot[R]{\sffamily\small\color{fautech}\thepage\ / \pageref*{LastPage}}
\renewcommand{\headrule}{\hbox to\headwidth{\color{fautech}\leaders\hrule height 0.4pt\hfill}}

\newtcolorbox{task}[1]{enhanced, sharp corners, frame hidden,
  colback=fautech!12, borderline west={2.5pt}{0pt}{faublue},
  left=9pt, right=9pt, top=6pt, bottom=6pt, after skip=12pt,
  before={\par\nopagebreak\vskip 6pt},
  fontupper=\sffamily\small,
  before upper={{\bfseries\color{faublue}Task~#1.}\enspace}}
\setcounter{MaxMatrixCols}{20}

\theoremstyle{remark}
\newtheorem{remark}{Remark}[section]

% ---------------------------------------------------------------
% Notation
% ---------------------------------------------------------------
\newcommand{\bigO}[1]{\mathcal{O}\!\left(#1\right)}
\newcommand{\Dint}{\textstyle\int_0^t D\,\mathrm{d}t'}

\begin{document}

\begin{titlepage}
  \noindent
  \includegraphics[height=2.0cm]{figures/fau-techfak.png}\hfill
  \raisebox{0.6cm}{\includegraphics[height=1.1cm]{figures/ww8.png}}

  \vspace{0.2cm}
  {\color{fautech}\rule{\textwidth}{0.6pt}}\\[0.1cm]
  {\sffamily\small\color{fautech}Institute of Materials Simulation (WW8)\hfill
    Winter semester 2026/27}

  \vspace*{3.4cm}
  {\sffamily\large\color{fautech}\textls[120]{MATERIALS SIMULATION PRACTICAL 6}\par}
  \vspace{0.5cm}
  {\sffamily\bfseries\color{faublue}\fontsize{30}{36}\selectfont
    Kinetic Monte Carlo\\[0.1em] Simulation of Diffusion\par}
  \vspace{0.5cm}
  {\sffamily\Large\color{faublue}Al in Cu Thin Films, Coarse Graining, and Au in Si by Kick-Out\par}
  \vspace{0.7cm}
  {\color{faublue}\rule{3cm}{2pt}}\par
  \vspace{0.4cm}
  {\sffamily\large Report, Part~I \textbullet{} Theory and Results\par}

  \vfill
  {\sffamily
  \renewcommand{\arraystretch}{1.35}
  \begin{tabular}{@{}>{\small\color{fautech}}l@{\hspace{1.2em}}l@{}}
    Author     & Ali Ghelichkhani \\
    E-mail     & \href{mailto:ali.ghelichkhani@fau.de}{ali.ghelichkhani@fau.de} \\
    Supervisor & Frank Wendler \\
    Module     & Foundations of Materials Simulation \\
    Institute  & Institute of Materials Simulation (WW8), FAU Erlangen-N\"urnberg \\
    Part~II    & Implementation, Jupyter notebook \texttt{kmc\_diffusion.ipynb}
  \end{tabular}\par}
\end{titlepage}

\section*{Abstract}
  We implement kinetic Monte Carlo (KMC) models for three diffusion problems of increasing complexity.
  For Al in a Cu/Cu--10\,\%Al thin-film sandwich on a 2D lattice, the jump rate follows from the
  macroscopic diffusion coefficient as $\Gamma = D/b^2$ per direction. Because every rate is
  proportional to $\Gamma(t)$, a time-dependent temperature enters only through the reduced time
  $\int\Gamma\,\mathrm{d}t$. Isothermal anneals at 500\,K and 600\,K reproduce the error-function solution
  to within the statistical noise. Heating from 300\,K to 873\,K in 1\,s gives a diffusion length of
  $\SI{8.3}{\nano\metre}$, $77\,\%$ of it accumulated in the last $\SI{100}{\milli\second}$. A
  coarse-grained description with cells of $\SI{1}{\micro\metre}$ and hop rates
  $(D/L^2)\,n_i(1 - n_j/n_{\max})$ treats $\SI{0.1}{\milli\metre}$ films heated over $10^6$\,s with about
  150 hops per diffusor, and in the scaling variable $x/\sqrt{\int D\,\mathrm dt}$ it coincides with the
  atomistic result. The blocked-site factor cancels in the net flux, so the collective diffusion
  coefficient stays concentration independent. For gold in a 1\,mm silicon wafer at 1050\,\textdegree C,
  a coarse-grained KMC of the kick-out reaction $\mathrm{Au}_i \rightleftharpoons \mathrm{Au}_s + I$
  with surface sinks for self-interstitials yields U-shaped substitutional profiles after 1000\,s. They
  are high at the gold-coated surface, two orders of magnitude lower in the bulk, and rise again
  towards the uncoated back surface, because gold incorporation is limited by the out-diffusion of
  self-interstitials. An independent solution of the rate equations confirms the KMC result.

\tableofcontents
\clearpage

% ===============================================================
\section{Introduction}\label{sec:intro}
% ===============================================================

Molecular dynamics resolves atomic vibrations and is therefore limited to nanoseconds
(Practical~4). Most of the processes that shape a microstructure --- diffusion, precipitation,
recrystallization --- are sequences of rare, thermally activated events separated by long waiting
times in which nothing of interest happens. Kinetic Monte Carlo~\citep{script6, voter2007} exploits
this separation of time scales. It does not follow the vibrations; it jumps directly from one event to
the next, with the waiting times and the choice of events drawn from the known transition rates. The
time step is then set by the events themselves and can reach seconds or days.

This report applies the method to three problems from the tutorial script: atomistic diffusion of Al
in a Cu thin film under a time-dependent temperature (\cref{sec:atomistic}), the same problem on much
larger length and time scales with a coarse-grained description (\cref{sec:coarse}), and the diffusion
of gold in silicon by the kick-out mechanism (\cref{sec:kickout}). The algorithm is summarized in
\cref{sec:kmc}; the code and all results are in Part~II.

% ===============================================================
\section{The kinetic Monte Carlo method}\label{sec:kmc}
% ===============================================================

Let the system be in state $i$, from which transitions $i\to j$ occur with rates $\nu_{ij}$. The
residence-time (BKL, $n$-fold way) algorithm~\citep{bortz1975, gillespie1977} proceeds as follows:
\begin{enumerate}\setlength{\itemsep}{1pt}
  \item compute the total rate $\Gamma = \sum_j\nu_{ij}$ and the cumulative probabilities
    $P_j = \sum_{j'\le j}\nu_{ij'}/\Gamma$;
  \item with a uniform random number $R_1\in(0,1]$ choose the event $j$ with $P_{j-1} < R_1 \le P_j$;
  \item with a second random number $R_2$ advance time by $t_{ij} = -\ln R_2/\Gamma$;
  \item execute the event, update the rates that changed, and repeat.
\end{enumerate}
The waiting time is exponentially distributed because, for independent Poisson processes, the time to
the first event of any of them is exponential with the total rate. The algorithm is exact for
time-independent rates and rejection-free. The physics lies entirely in the catalogue of events and
their rates; the cost lies in finding the event, which is why it pays to split the system into
independent parts and to update the rate sums incrementally.

\paragraph{Time-dependent rates.}
If all rates share a common time dependence, $\nu_{ij}(t) = g(t)\,w_{ij}$, the dynamics is a
time-homogeneous process in the reduced time $\tau = \int_0^t g(t')\,\mathrm{d}t'$: one simulates with
the constant rates $w_{ij}$ in $\tau$ and maps back to physical time by inverting $\tau(t)$. This is
exact for any heating schedule, and it is the situation of \cref{sec:atomistic,sec:coarse}, where every
jump rate is proportional to the same Arrhenius factor.

% ===============================================================
\section{Atomistic KMC: Al in a Cu thin film (Problem 2.1)}\label{sec:atomistic}
% ===============================================================

\begin{task}{2.1}
  A $\SI{100}{\nano\metre}$ Cu film is in contact with a $\SI{100}{\nano\metre}$ Cu--10\,\%Al film. Al
  sits on substitutional sites; $D = 1.49\cdot10^{-7}\exp(-137.1\,\mathrm{kJ}/RT)$\,m$^2$/s and
  $b = \SI{2.54}{\angstrom}$. Starting at 300\,K, the temperature rises linearly to 873\,K within 1\,s.
  Evaluate the Al profile over this time, after verifying the implementation with isothermal anneals at
  500\,K and 600\,K against the error-function solution. Use a 2D lattice and remember that an Al atom
  cannot jump onto a site occupied by another Al atom.
\end{task}

\subsection{Model and rates}\label{sec:atmodel}

\paragraph{Geometry.}
The problem is homogeneous in the plane of the film, so a 2D square lattice suffices: $x$ is the film
normal with $2\times394$ lattice planes (two times $\SI{100}{\nano\metre}$), $y$ a periodic direction of
100 sites representing the film plane. The outer faces are reflecting. Initially $10\,\%$ of the sites
of the right half are occupied by Al at random.

\paragraph{Jump rate.}
On a lattice of coordination $z$ in $d$ dimensions with jump rate $\Gamma$ per direction, the
Einstein relation gives $D = z\Gamma b^2/(2d)$. For the 2D square lattice, $z = 4$ and $d = 2$, hence
\begin{equation}\label{eq:gamma}
  \Gamma(T) = \frac{D(T)}{b^2} = \frac{D_0}{b^2}\exp\!\left(-\frac{Q}{RT}\right),
\end{equation}
which ranges from $3\cdot10^{-12}\,\mathrm{s^{-1}}$ at 300\,K to $1.4\cdot10^{4}\,\mathrm{s^{-1}}$ at
873\,K. Physically, Al moves by exchange with vacancies; the macroscopic $D$ is the dilute-limit
impurity diffusion coefficient measured by Laik et al.~\citep{laik2008}, and \cref{eq:gamma} is the
rate of an effective Al--Cu exchange that reproduces it.

\paragraph{Blocked sites.}
A jump onto a site occupied by another Al atom is forbidden. All allowed jumps have the same rate, so
the BKL selection amounts to choosing uniformly among the allowed (atom, direction) pairs. It is
simpler, and exactly equivalent, to choose uniformly among all $4N$ pairs and reject the forbidden ones
as null events with time advancing by $-\ln R_2/(4N\Gamma)$; the rejected attempts then carry precisely
the waiting time that the rejection-free algorithm would assign to the allowed ones. This model is the
symmetric exclusion process, for which the expected occupation $\langle n_k\rangle$ of a site obeys
the discrete diffusion equation \emph{exactly}, because the blocked contributions cancel pairwise. The
collective (chemical) diffusion coefficient is therefore $D$ at any concentration, and the error
function is the exact reference solution. Only the tracer diffusion of an individual atom is slowed
down by blocking.

\subsection{Verification at constant temperature}\label{sec:isothermal}

For two semi-infinite regions with $c = 0$ and $c = c_0$ the solution of Fick's second law is
\begin{equation}\label{eq:erf}
  c(x,t) = \frac{c_0}{2}\left[1 + \mathrm{erf}\!\left(\frac{x - x_0}{2\sqrt{Dt}}\right)\right],
\end{equation}
valid as long as the diffusion length is small against the film thickness. Part~II anneals at 500\,K
for $3.5\cdot10^{4}$ and $1.4\cdot10^{5}$\,s and at 600\,K for 145 and 578\,s, which gives
$\sqrt{Dt} = 5$ and $\SI{10}{\nano\metre}$ in both cases. Averaged over the film plane, over four runs
and over $\SI{5}{\nano\metre}$ slabs, the KMC profiles deviate from \cref{eq:erf} by an rms of
$0.0024$--$0.0028$, which is the binomial noise of the slab averages. The two temperatures give
statistically identical profiles at the same $Dt$, which illustrates the reduced-time argument of
\cref{sec:kmc}: temperature only sets the clock.

\subsection{Heating to 873\,K in one second}\label{sec:ramp}

For $T(t) = 300\,\mathrm{K} + 573\,\mathrm{K}\cdot t/\mathrm{s}$ the profile is \cref{eq:erf} with
$Dt$ replaced by $\int_0^t D(T(t'))\,\mathrm{d}t'$. \Cref{tab:ramp} lists the values, and Part~II shows
that the KMC profiles agree with them at all times.

\begin{table}[ht]
  \centering
  \caption{Accumulated diffusion during heating from 300\,K to 873\,K in 1\,s (Part~II).}
  \label{tab:ramp}
  \begin{tabular}{rrrrr}
    \toprule
    $t$ [s] & $T$ [K] & $\int_0^t D\,\mathrm{d}t'$ [m$^2$] & $\sqrt{\int D\,\mathrm{d}t'}$ [nm]
      & jumps per Al atom \\
    \midrule
    0.50 & 586.5 & $3.1\cdot10^{-21}$ & 0.06 & 0.2 \\
    0.80 & 758.4 & $3.0\cdot10^{-18}$ & 1.73 & 187 \\
    0.90 & 815.7 & $1.6\cdot10^{-17}$ & 3.99 & 988 \\
    0.95 & 844.4 & $3.4\cdot10^{-17}$ & 5.81 & 2096 \\
    1.00 & 873.0 & $6.8\cdot10^{-17}$ & 8.27 & 4240 \\
    \bottomrule
  \end{tabular}
\end{table}

The Arrhenius dependence makes the history extremely uneven. During the first half second, below
590\,K, an Al atom makes on average less than one jump; $77\,\%$ of the total diffusion takes place in
the last $\SI{100}{\milli\second}$, above 816\,K. A useful approximation follows from expanding
$1/T$ about the final temperature:
$\int D\,\mathrm{d}t \approx D(T_{\max})\,(RT_{\max}^2/Q)/\dot T$, i.e.\ the ramp acts like an isothermal
anneal at $T_{\max}$ for the time $RT_{\max}^2/(Q\dot T) \approx \SI{81}{\milli\second}$. At the end of
the ramp the interdiffusion zone is only $\SI{8}{\nano\metre}$ wide: the sandwich is essentially intact.

% ===============================================================
\section{Coarse-grained KMC (Problem 2.2)}\label{sec:coarse}
% ===============================================================

\begin{task}{2.2}
  Treat the same problem with layer thicknesses of $2\times\SI{0.1}{\milli\metre}$ and a heating time of
  $10^6$\,s. Instead of atomic jumps use transitions of diffusors between cells of size $L$, with rates
  chosen so that the overall diffusion coefficient is the same as in the atomistic description. Deal
  with the concentration dependence caused by blocked sites, project the process on a single
  coordinate, and compare with Problem~2.1 for the same number of diffusors.
\end{task}

\paragraph{Why coarse grain.}
The atomistic model would need $4\int D\,\mathrm{d}t/b^2 \approx 5\cdot10^{9}$ jumps per atom for this
problem and a lattice of $\sim10^{6}$ planes. Lumping many sites into a cell removes both obstacles.

\paragraph{Rates.}
Project onto $x$ and divide it into cells of length $L$, each containing $n_{\max}$ lattice sites. A
random walk between neighbouring cells with hop rate $w$ per direction has $D = wL^2$, so matching the
macroscopic coefficient in the dilute limit requires $w = D/L^2$. For a cell with $n_i$ diffusors next
to a cell with $n_j$, the total rate for one diffusor to move from $i$ to $j$ is
\begin{equation}\label{eq:cgrate}
  W_{i\to j} = \frac{D(T(t))}{L^2}\,n_i\left(1 - \frac{n_j}{n_{\max}}\right),
\end{equation}
where the factor in parentheses is the probability that the target site is free --- the mean-field
version of the blocked-site rule.

\paragraph{The blocked-site problem.}
The expected net flux between two cells is
\begin{equation}\label{eq:cgflux}
  W_{i\to j} - W_{j\to i}
  = \frac{D}{L^2}\left[n_i\left(1 - \frac{n_j}{n_{\max}}\right) - n_j\left(1 - \frac{n_i}{n_{\max}}\right)\right]
  = \frac{D}{L^2}\,(n_i - n_j),
\end{equation}
so the blocking factors cancel and the collective diffusion coefficient is concentration independent,
exactly as in the atomistic exclusion process of \cref{sec:atmodel}. The factor must nevertheless be
included: it keeps every cell below its capacity and makes the coarse-grained model the projection of
the atomistic one, in which each individual diffusor \emph{is} slowed down. Had the measured $D$
depended on concentration, as real Cu--Al interdiffusion coefficients do~\citep{laik2008}, it would be
evaluated at the mean concentration of the two cells, $D(\bar c_{ij})$.

\paragraph{Results.}
We use cells of $L = \SI{1}{\micro\metre}$, 200 cells across the sandwich and the same 3940 diffusors
as in Problem~2.1, which gives $n_{\max} = 394$ and $10\,\%$ occupation in the Al-rich layer. All rates
are proportional to $D(t)$, so the simulation runs in reduced time $\int D\,\mathrm{d}t/L^2$ with the
BKL algorithm. Heating over $10^6$\,s accumulates $\sqrt{\int D\,\mathrm{d}t} = \SI{8.3}{\micro\metre}$
--- the same number as in Problem~2.1 in micrometres instead of nanometres, because the ramp is
$10^6$ times slower and the diffusion length scales with the square root of the time. A diffusor then
makes only about 150 hops. The profiles agree with the error function to an rms of
$0.003$--$0.004$, with and without the blocking factor, as \cref{eq:cgflux} predicts. In the scaling
variable $(x-x_0)/\sqrt{\int D\,\mathrm{d}t}$ the atomistic and the coarse-grained profiles fall on the
same curve: both are realizations of the same diffusion equation, sampled on scales that differ by a
factor of a thousand in length and a million in time.

% ===============================================================
\section{Au in Si by the kick-out mechanism (Problem 2.3)}\label{sec:kickout}
% ===============================================================

\begin{task}{2.3}
  Gold diffuses into a $\SI{1}{\milli\metre}$ Si wafer at 1050\,\textdegree C. Interstitial gold
  $\mathrm{Au}_i$ is fast, substitutional gold $\mathrm{Au}_s$ immobile, and they interchange by
  $\mathrm{Au}_i + \mathrm{Si} \rightleftharpoons \mathrm{Au}_s + I$. At the left surface
  $\mathrm{Au}_s$ is kept at its solubility $5.38\cdot10^{-7}$; the surfaces are ideal sinks for
  self-interstitials $I$; $D_{\mathrm{Au}} \approx D_I = 10^{-9}$\,m$^2$/s; $a = \SI{5.4}{\angstrom}$;
  the kick-out rates are 600\,s$^{-1}$ for $\mathrm{Au}_i$ and 10\,s$^{-1}$ for $I$. Set up a
  coarse-grained simulation and evaluate the profiles after 1000\,s. Explain their shape.
\end{task}

\subsection{Model}\label{sec:komodel}

The wafer is divided into 100 cells of $\Delta x = \SI{10}{\micro\metre}$; each cell holds $n_i$
interstitial gold atoms, $n_s$ substitutional gold atoms and $n_I$ self-interstitials. With
$C_0 = 8/a^3 = 5.1\cdot10^{28}\,\mathrm{m^{-3}}$ silicon sites, a cell cross-section of
$(\SI{43}{\nano\metre})^2$ makes the solubility correspond to $N_s^{\mathrm{eq}} = 500$ substitutional
atoms per cell. The events are:
\begin{itemize}\setlength{\itemsep}{2pt}
  \item hops of $\mathrm{Au}_i$ and $I$ to a neighbouring cell, $D/\Delta x^2 = 10\,\mathrm{s^{-1}}$ per
    particle and direction;
  \item kick-out of a Si atom by an interstitial gold atom, $\mathrm{Au}_i \to \mathrm{Au}_s + I$, with
    $r_f = k_f\,n_i$ and $k_f = \SI{600}{\per\second}$, since practically every neighbour of an
    $\mathrm{Au}_i$ is a Si atom;
  \item the reverse reaction $\mathrm{Au}_s + I \to \mathrm{Au}_i$, in which a self-interstitial
    kicks out a substitutional gold atom, with
    \begin{equation}\label{eq:reverse}
      r_r = k_r\,n_s\,\frac{n_I}{N_I^{\mathrm{eq}}},\qquad k_r = \SI{10}{\per\second}.
    \end{equation}
\end{itemize}
The $10\,\mathrm{s^{-1}}$ is the rate at which a substitutional gold atom is kicked out when the
self-interstitials are at equilibrium, and the rate scales with the local $I$ concentration. This is
exactly the kick-out model of Lerch and Stolwijk~\citep{lerch1998}: at 1050\,\textdegree C they give
$k_2C_I^{\mathrm{eq}} = \SI{10}{\per\second}$, $C_i^{\mathrm{eq}}/C_s^{\mathrm{eq}} = 0.0166$ and
$C_I^{\mathrm{eq}}/C_s^{\mathrm{eq}} = 6.68\cdot10^{-5}$, so that the mass-action relation
$k_1 = k_2C_s^{\mathrm{eq}}C_I^{\mathrm{eq}}/C_i^{\mathrm{eq}}$ yields $k_1 \approx \SI{600}{\per\second}$
--- the two rates of the script. In local equilibrium $r_f = r_r$ becomes
\begin{equation}\label{eq:localeq}
  \frac{C_i}{C_i^{\mathrm{eq}}} = \frac{C_s}{C_s^{\mathrm{eq}}}\,\frac{C_I}{C_I^{\mathrm{eq}}}.
\end{equation}
Boundary conditions: keeping $\mathrm{Au}_s$ at its solubility at the coated surface means, through
\cref{eq:localeq}, keeping $\mathrm{Au}_i$ at its solubility
$C_i^{\mathrm{eq}} = (k_r/k_f)\,C_s^{\mathrm{eq}}$ there; it is implemented as a ghost cell that feeds
$\mathrm{Au}_i$ into the first cell and absorbs atoms hopping back. Both surfaces annihilate $I$, and
$\mathrm{Au}_i$ cannot leave through the uncoated back surface.

The kick-out is fast compared with transport: an interstitial gold atom lives $1/k_f = \SI{1.7}{\milli\second}$
and travels only $\sqrt{D/k_f} = \SI{1.3}{\micro\metre}$ before it is incorporated, unless the reverse
reaction frees it again. The simulation therefore spends almost all events on the reaction pair, about
$4\cdot10^8$ events per run, which is why the event loop is compiled. The deterministic rate
equations of the same model are solved as an independent check.

\subsection{Results: the ``weird'' profiles}\label{sec:koresults}

After 1000\,s the substitutional gold profile (Part~II) is not an erfc decaying from the coated surface
but U-shaped. $C_s/C_s^{\mathrm{eq}}$ is $0.47$ in the first cell, falls to a minimum of $0.013$ at a
depth of $\SI{0.7}{\milli\metre}$, and rises again to $0.13$ at the uncoated back surface, which has no
gold source. The self-interstitials show the inverse shape, a broad maximum of up to about $30\,C_I^{\mathrm{eq}}$
in the bulk, and the interstitial gold decreases only gradually through the wafer.
The KMC and the rate equations agree within the noise of the KMC. The explanation:
\begin{enumerate}\setlength{\itemsep}{2pt}
  \item Interstitial gold is mobile and penetrates the whole wafer within the annealing time
    ($\sqrt{Dt} = \SI{1}{\milli\metre}$), but it is rare: its solubility is only $1.7\,\%$ of that of
    substitutional gold.
  \item Each incorporation creates a self-interstitial. A dislocation-free crystal has no internal
    sinks, so in the bulk the $I$ accumulate until the reverse reaction balances the forward one. With
    $C_I \gg C_I^{\mathrm{eq}}$, \cref{eq:localeq} allows only a small $C_s$: incorporation in the bulk
    is choked and proceeds only as fast as the $I$ can diffuse out.
  \item Near both surfaces the $I$ escape quickly, $C_I$ stays close to equilibrium, and the gold
    arriving as $\mathrm{Au}_i$ is incorporated freely. The gold enters at the front, but it is
    incorporated wherever the self-interstitials can leave, and that includes the back surface.
\end{enumerate}
The rate of gold incorporation is therefore controlled not by the diffusion of gold but by the
out-diffusion of the self-interstitials it produces, which is the key result of the kick-out analysis
of Au, Pt and Zn in silicon~\citep{gosele1980, lerch1998}. The same mechanism explains why the surface
concentration approaches the solubility only slowly, as Lerch and Stolwijk observed experimentally, and
why the gold profiles in dislocated silicon, where dislocations act as bulk sinks for $I$, look
completely different.

% ===============================================================
\section{Conclusion}\label{sec:conclusion}
% ===============================================================

Kinetic Monte Carlo reduces the dynamics of a system to a catalogue of events and their rates, and
everything else follows from the residence-time algorithm. For substitutional diffusion on a lattice
the rate per jump is $D/b^2$, blocking turns the model into an exclusion process whose collective
diffusion is still exactly Fickian, and a time-dependent temperature only rescales time. With these
three ingredients the atomistic simulation of the Cu/Cu--Al sandwich reproduces the error-function
profiles at 500\,K and 600\,K and shows that a 1\,s heating ramp to 873\,K interdiffuses only
$\SI{8}{\nano\metre}$, almost all of it in the final $\SI{100}{\milli\second}$. Coarse graining with hop
rates $(D/L^2)\,n_i(1-n_j/n_{\max})$ preserves the diffusion coefficient by construction, handles films
a thousand times thicker heated a million times longer with about 150 hops per diffusor, and gives the
same profile in the scaling variable. For gold in silicon the coarse-grained KMC of the kick-out
reaction produces the characteristic U-shaped profiles with a deep minimum in the bulk and a second
maximum at the uncoated surface: gold incorporation is limited by the escape of the self-interstitials
that the reaction creates, not by the transport of gold itself.

\begin{thebibliography}{99}

\bibitem{script6}
Institute of Materials Simulation (WW8),
``Kinetic Monte Carlo simulation,''
Materials Simulation Practical, tutorial script,
Friedrich-Alexander-Universit\"at Erlangen-N\"urnberg.

\bibitem{voter2007}
A.~F.~Voter,
``Introduction to the kinetic Monte Carlo method,''
in \textit{Radiation Effects in Solids}, K.~E.~Sickafus, E.~A.~Kotomin and B.~P.~Uberuaga (Eds.),
Springer, Dordrecht, 2007, pp.~1--23.

\bibitem{bortz1975}
A.~B.~Bortz, M.~H.~Kalos and J.~L.~Lebowitz,
``A new algorithm for Monte Carlo simulation of Ising spin systems,''
\textit{Journal of Computational Physics},
vol.~17, no.~1, pp.~10--18, 1975.

\bibitem{gillespie1977}
D.~T.~Gillespie,
``Exact stochastic simulation of coupled chemical reactions,''
\textit{The Journal of Physical Chemistry},
vol.~81, no.~25, pp.~2340--2361, 1977.

\bibitem{laik2008}
A.~Laik, K.~Bhanumurthy and G.~B.~Kale,
``Diffusion in Cu(Al) solid solution,''
\textit{Defect and Diffusion Forum},
vol.~279, pp.~63--69, 2008.

\bibitem{lerch1998}
W.~Lerch and N.~A.~Stolwijk,
``Diffusion of gold in silicon during rapid thermal annealing: Effectiveness of the surface as a sink
for self-interstitials,''
\textit{Journal of Applied Physics},
vol.~83, no.~3, pp.~1312--1320, 1998.

\bibitem{gosele1980}
U.~G\"osele, W.~Frank and A.~Seeger,
``Mechanism and kinetics of the diffusion of gold in silicon,''
\textit{Applied Physics},
vol.~23, no.~4, pp.~361--368, 1980.

\end{thebibliography}

\end{document}
